% chktex-file 1
% chktex-file 2
% chktex-file 3
% chktex-file 8
% chktex-file 11
% chktex-file 12
% chktex-file 24
% chktex-file 26
% chktex-file 36
% chktex-file 44

\DefineTable{tab0301}{
  \begingroup
  \small
  \addtolength{\abovecaptionskip}{-0.2cm}
  \setlength{\tabcolsep}{3pt}
  \setlength{\dashlinedash}{1.5pt}
  \setlength{\dashlinegap}{1.5pt}
  \renewcommand{\arraystretch}{1.10}
  \edef\symboltablearrayrulewidth{\the\arrayrulewidth}
  \newcommand{\symboltablerule}[1]{%
    \noalign{\global\arrayrulewidth=#1\relax}%
    \hline
    \noalign{\global\arrayrulewidth=\symboltablearrayrulewidth\relax}%
  }
  \begin{longtable}{
    >{\centering\arraybackslash}p{0.160\textwidth}
    >{\centering\arraybackslash}p{0.300\textwidth}:
    >{\centering\arraybackslash}p{0.160\textwidth}
    >{\centering\arraybackslash}p{0.300\textwidth}}
    \caption{第三章符号汇总}\label{tab0301}                                                                                                                                                                                  \\
    \symboltablerule{\heavyrulewidth}
    \multicolumn{4}{c}{\cellcolor{gray!30}\textbf{时间}}                                                                                                                                                                \\*
    \multicolumn{1}{c}{\cellcolor{gray!30}\textbf{符号}} & \multicolumn{1}{c}{\cellcolor{gray!30}\textbf{含义}} & \multicolumn{1}{c}{\cellcolor{gray!30}\textbf{符号}} & \multicolumn{1}{c}{\cellcolor{gray!30}\textbf{含义}} \\*
    \symboltablerule{\lightrulewidth}
    $T_x$                                              & 周期                                                 & $f_x$                                              & 频率                                                 \\
    $D$                                                & 有效占空比                                              & $D'$                                               & 互补占空比                                              \\
    $D_{\mathrm{c}}$                                   & 控制信号占空比                                            & $c(t)$                                             & 开关控制信号                                             \\
    $t_x$                                              & 器件特征时间                                             & $\Delta t_x$                                       & 开关过渡时间                                             \\
    \symboltablerule{\heavyrulewidth}
    \multicolumn{4}{c}{\cellcolor{gray!30}\textbf{电压}}                                                                                                                                                                \\*
    \multicolumn{1}{c}{\cellcolor{gray!30}\textbf{符号}} & \multicolumn{1}{c}{\cellcolor{gray!30}\textbf{含义}} & \multicolumn{1}{c}{\cellcolor{gray!30}\textbf{符号}} & \multicolumn{1}{c}{\cellcolor{gray!30}\textbf{含义}} \\*
    \symboltablerule{\lightrulewidth}
    $V_x$                                              & 直流分量或平均值                                           & $v_x(t)$                                           & 瞬时量                                                \\
    $V_{x(y)}$                                         & 端口 $x$ 在条件 $y$ 下的电压                                & $V_{xy}$                                           & 端子 $x$ 对 $y$ 的电压                                   \\
    \symboltablerule{\heavyrulewidth}
    \multicolumn{4}{c}{\cellcolor{gray!30}\textbf{电流}}                                                                                                                                                                \\*
    \multicolumn{1}{c}{\cellcolor{gray!30}\textbf{符号}} & \multicolumn{1}{c}{\cellcolor{gray!30}\textbf{含义}} & \multicolumn{1}{c}{\cellcolor{gray!30}\textbf{符号}} & \multicolumn{1}{c}{\cellcolor{gray!30}\textbf{含义}} \\*
    \symboltablerule{\lightrulewidth}
    $I_x$                                              & 直流分量或平均值                                           & $i_x(t)$                                           & 瞬时量                                                \\
    $I_{x,y}$                                          & 电流统计量                                              &                                                    &                                                    \\
    \symboltablerule{\heavyrulewidth}
    \multicolumn{4}{c}{\cellcolor{gray!30}\textbf{电荷与器件参数}}                                                                                                                                                           \\*
    \multicolumn{1}{c}{\cellcolor{gray!30}\textbf{符号}} & \multicolumn{1}{c}{\cellcolor{gray!30}\textbf{含义}} & \multicolumn{1}{c}{\cellcolor{gray!30}\textbf{符号}} & \multicolumn{1}{c}{\cellcolor{gray!30}\textbf{含义}} \\*
    \symboltablerule{\lightrulewidth}
    $Q_x$                                              & 电荷                                                 & $R_x$                                              & 电阻                                                 \\
    $C_{xy}$                                           & 端子 $x$、$y$ 间的电容                                    &                                                    &                                                    \\
    \symboltablerule{\heavyrulewidth}
    \multicolumn{4}{c}{\cellcolor{gray!30}\textbf{功率、能量与效率}}                                                                                                                                                          \\*
    \multicolumn{1}{c}{\cellcolor{gray!30}\textbf{符号}} & \multicolumn{1}{c}{\cellcolor{gray!30}\textbf{含义}} & \multicolumn{1}{c}{\cellcolor{gray!30}\textbf{符号}} & \multicolumn{1}{c}{\cellcolor{gray!30}\textbf{含义}} \\*
    \symboltablerule{\lightrulewidth}
    $P_x$                                              & 平均功率                                               & $p_x(t)$                                           & 瞬时功率                                               \\
    $E_x$                                              & 单次开关损耗能量                                           & $\eta$                                             & 变换器效率                                              \\
    $M(D)$                                             & 直流电压变换比                                            &                                                    &                                                    \\
    \symboltablerule{\heavyrulewidth}
  \end{longtable}

  \endgroup
}

\DefineTable{tab0302}{
  \begingroup
  \small
  \addtolength{\abovecaptionskip}{-0.2cm}
  \setlength{\tabcolsep}{3pt}
  \setlength{\dashlinedash}{1.5pt}
  \setlength{\dashlinegap}{1.5pt}
  \renewcommand{\arraystretch}{1.10}
  \edef\symboltablearrayrulewidth{\the\arrayrulewidth}
  \newcommand{\symboltablerule}[1]{%
    \noalign{\global\arrayrulewidth=#1\relax}%
    \hline
    \noalign{\global\arrayrulewidth=\symboltablearrayrulewidth\relax}%
  }
  \begin{longtable}{
    >{\centering\arraybackslash}p{0.160\textwidth}
    >{\centering\arraybackslash}p{0.300\textwidth}:
    >{\centering\arraybackslash}p{0.160\textwidth}
    >{\centering\arraybackslash}p{0.300\textwidth}}
    \caption{第三章符号通用标记汇总}\label{tab0302}                                                                                                                                                                              \\
    \symboltablerule{\heavyrulewidth}
    \multicolumn{4}{c}{\cellcolor{gray!30}\textbf{数学与统计标记}}                                                                                                                                                           \\*
    \multicolumn{1}{c}{\cellcolor{gray!30}\textbf{符号}} & \multicolumn{1}{c}{\cellcolor{gray!30}\textbf{含义}} & \multicolumn{1}{c}{\cellcolor{gray!30}\textbf{符号}} & \multicolumn{1}{c}{\cellcolor{gray!30}\textbf{含义}} \\*
    \symboltablerule{\lightrulewidth}
    $x_{y,\mathrm{av}}$                                & 平均值                                                & $x_{y,\mathrm{rms}}$                               & 有效值                                                \\
    $\Delta x_{y,\mathrm{pp}}$                         & 峰峰值                                                & $\langle x_y\rangle$                               & 单周期内均值                                             \\
    $x_{1},\;x_{2},\;\cdots$                           & 量、元器件或时刻的序号                                        &                                                    &                                                    \\
    \symboltablerule{\heavyrulewidth}
    \multicolumn{4}{c}{\cellcolor{gray!30}\textbf{电路位置与对象标记}}                                                                                                                                                         \\*
    \multicolumn{1}{c}{\cellcolor{gray!30}\textbf{符号}} & \multicolumn{1}{c}{\cellcolor{gray!30}\textbf{含义}} & \multicolumn{1}{c}{\cellcolor{gray!30}\textbf{符号}} & \multicolumn{1}{c}{\cellcolor{gray!30}\textbf{含义}} \\*
    \symboltablerule{\lightrulewidth}
    $x_{\mathrm{in}}$                                  & 输入                                                 & $x_{\mathrm{out}}$                                 & 输出                                                 \\
    $x_{\mathrm L}$                                    & 电感                                                 & $x_{\mathrm C}$                                    & 电容                                                 \\
    $x_{\mathrm D}$                                    & 二极管                                                & $x_{\mathrm Q}$                                    & 开关管                                                \\
    $x_{\mathrm s}$                                    & 开关                                                 &                                                    &                                                    \\
    \symboltablerule{\heavyrulewidth}
    \multicolumn{4}{c}{\cellcolor{gray!30}\textbf{器件端口标记}}                                                                                                                                                            \\*
    \multicolumn{1}{c}{\cellcolor{gray!30}\textbf{符号}} & \multicolumn{1}{c}{\cellcolor{gray!30}\textbf{含义}} & \multicolumn{1}{c}{\cellcolor{gray!30}\textbf{符号}} & \multicolumn{1}{c}{\cellcolor{gray!30}\textbf{含义}} \\*
    \symboltablerule{\lightrulewidth}
    $x_{\mathrm G}$                                    & MOSFET、IGBT 的栅极                                    & $x_{\mathrm G}$                                    & SCR 的门极                                            \\
    $x_{\mathrm D}$                                    & MOSFET 的漏极                                         & $x_{\mathrm S}$                                    & MOSFET 的源极                                         \\
    $x_{\mathrm B}$                                    & BJT 的基极                                            & $x_{\mathrm C}$                                    & BJT、IGBT 的集电极                                      \\
    $x_{\mathrm E}$                                    & BJT、IGBT 的发射极                                      & $x_{\mathrm A},\;x_{\mathrm K}$                    & SCR 的阳极、阴极                                         \\
    $x_{\mathrm{GS}}$                                  & 栅极、源极之间                                            & $x_{\mathrm{DS}}$                                  & 漏极、源极之间                                            \\
    $x_{\mathrm{BE}}$                                  & 基极、发射极之间                                           & $x_{\mathrm{CE}}$                                  & 集电极、发射极之间                                          \\
    $x_{\mathrm{GE}}$                                  & 栅极、发射极之间                                           & $x_{\mathrm{AK}}$                                  & 阳极、阴极之间                                            \\
    \symboltablerule{\heavyrulewidth}
    \multicolumn{4}{c}{\cellcolor{gray!30}\textbf{状态、特性与损耗标记}}                                                                                                                                                        \\*
    \multicolumn{1}{c}{\cellcolor{gray!30}\textbf{符号}} & \multicolumn{1}{c}{\cellcolor{gray!30}\textbf{含义}} & \multicolumn{1}{c}{\cellcolor{gray!30}\textbf{符号}} & \multicolumn{1}{c}{\cellcolor{gray!30}\textbf{含义}} \\*
    \symboltablerule{\lightrulewidth}
    $x_{\mathrm{on}}$                                  & 开通或通态                                              & $x_{\mathrm{off}}$                                 & 关断或断态                                              \\
    $x_{\mathrm d}$                                    & 延迟                                                 & $x_{\mathrm r}$                                    & 上升                                                 \\
    $x_{\mathrm f}$                                    & 下降                                                 & $x_{\mathrm q}$                                    & SCR 关断恢复                                           \\
    $x_{\mathrm F}$                                    & 二极管正向导通                                            & $x_{\mathrm{BO}}$                                  & 正向转折                                               \\
    $x_{\mathrm L}$                                    & SCR 锁定                                             & $x_{\mathrm H}$                                    & SCR 维持                                             \\
    $x_{y(\mathrm{sat})}$                              & 饱和状态                                               & $x_{y(\mathrm{th})}$                               & 阈值条件                                               \\
    $x_{\mathrm{rr}}$                                  & 反向恢复                                               & $x_{\mathrm{stg}}$                                 & 存储阶段                                               \\
    $x_{\mathrm{cond}}$                                & 通态损耗                                               & $x_{\mathrm{sw}}$                                  & 开关损耗                                               \\
    $x_{\mathrm{drv}}$                                 & 栅极驱动损耗                                             & $x_{\mathrm{loss}}$                                & 总损耗                                                \\
    \symboltablerule{\heavyrulewidth}
  \end{longtable}

  \endgroup
}
